\section*{अध्याय \ref{ch:g9:heredity-and-traits} --- आनुवंशिकता और लक्षण}

\begin{solution}{exo:g9:heredity-and-traits:1}
\omterm{def:g9:heredity-and-traits:traits}{लक्षण} किसी व्यक्ति की कोई भी बताई जा सकने वाली ख़ासियत है; और \omterm{def:g9:heredity-and-traits:heredity}{आनुवंशिकता}
जनकों से संतान तक, जनन के ज़रिए, \omterm{def:g9:heredity-and-traits:traits}{लक्षणों} का पहुँचना। छाँटने की कसौटी: जो
दो \omterm{def:g8:sexual-reproduction:gametes}{युग्मक} ढो नहीं सकते, परिवार वह पहुँचा नहीं सकते।
\end{solution}

\begin{solution}{exo:g9:heredity-and-traits:2}
रक्त-समूह: आनुवंशिक। छिदा हुआ \omterm{def:g1:the-five-senses:sense}{कान}: हासिल किया हुआ। झाइयाँ: एक आनुवंशिक
रुझान, और धूप में जिया हुआ (दोनों कारीगरियाँ)। फ़र्राटेदार फ़्रांसीसी: हासिल
की हुई। \omterm{def:g1:the-five-senses:sense}{कान} की लौ की शक्ल: आनुवंशिक।
\end{solution}

\begin{solution}{exo:g9:heredity-and-traits:3}
परिवार का ढर्रा (जनन की लकीरों के साथ लौटता है, या जीवन के पीछे चलता
है); शुरू से हाज़िर होना (\omterm{def:g5:flower-to-fruit:steps}{निषेचन} पर तय, या बाद में बना); और पहुँचना
(संतान तक जाता है, या अपने हासिल करने वाले के साथ मर जाता है)।
\end{solution}

\begin{solution}{exo:g9:heredity-and-traits:4}
क्योंकि सधी हुई \omterm{def:g3:bones-and-muscles:muscle}{पेशी} हासिल की हुई है --- इस्तेमाल से बनी, \omterm{def:g8:sexual-reproduction:gametes}{युग्मकों} की
ढोई हुई नहीं। यह दिखाता है कि किसी जीवन ने किसी शरीर में जो जोड़ा वह उसमें
लिखा ही नहीं जाता जो जनन पहुँचाता है।
\end{solution}

\begin{solution}{exo:g9:heredity-and-traits:5}
वह पहुँचा सकता है जो वह ख़ुद दिखाता नहीं: उस \omterm{def:g9:heredity-and-traits:traits}{लक्षण} का निर्धारक उसमें से
छिपा हुआ गुज़र गया, और एक पीढ़ी बाद ऊपर आ गया।
\end{solution}

\begin{solution}{exo:g9:heredity-and-traits:6}
क्योंकि ख़ुद \omterm{def:g9:heredity-and-traits:traits}{लक्षण} (कोई ठुड्डी, कोई रक्त-समूह) दो अकेली \omterm{def:g5:first-look-at-cells:cell}{कोशिकाओं} में सवार
हो ही नहीं सकते --- जो सवार होता है वह उन्हें बनाने के निर्देश ही होंगे।
\end{solution}

\begin{solution}{exo:g9:heredity-and-traits:7}
विरासत में मिला: वह लंबाई की सीमा, जिसका इशारा यह चौखट देती है ---
पोते को परिवार के निर्धारक मिले। और जिया हुआ: समृद्ध खुराक ने उस सीमा को
और आगे तक भरा --- यानी विरासत के ऊपर \omterm{def:g6:exploring-our-environment:components}{पर्यावरण} का जोड़।
\end{solution}

\begin{solution}{exo:g9:heredity-and-traits:8}
हर संतान को हर जनक के निर्धारकों का एक अलग चुनाव मिलता है --- फेंटाई हर
\omterm{def:g5:flower-to-fruit:steps}{निषेचन} पर ताज़े मेल बाँटती है। जबकि समरूप जुड़वाँ एक ही \omterm{def:g5:flower-to-fruit:steps}{निषेचन} के
\emph{बाद} बँटते हैं: एक बँटवारा, दो बनावटें --- इसलिए उनके निर्धारकों
के सेट मिलते हैं।
\end{solution}

\begin{solution}{exo:g9:heredity-and-traits:9}
(क) वंश-वृक्ष के पीछे चलता है, पेशे या शहर की परवाह किए बिना अपनी
पारिवारिक उम्र पर आता है, और \omterm{def:g8:sexual-reproduction:gametes}{युग्मकों} से ढोया जाना मुमकिन है: आनुवंशिक।
(ख) प्रशिक्षण से बना, खेल के पीछे चलता है जनकों के नहीं, और कोई \omterm{def:g8:sexual-reproduction:gametes}{युग्मक}
दमख़म ढोता ही नहीं: हासिल किया हुआ --- पर गठन तथा \omterm{def:g4:heart-and-blood:heart}{दिल} की विरासत में मिली
एक सीमा पर (\cref{prop:g7:organs-at-work:training} वही बनाती है जिसका
ख़ाका \omterm{def:g9:heredity-and-traits:heredity}{आनुवंशिकता} खींचती है)।
\end{solution}

\begin{solution}{exo:g9:heredity-and-traits:10}
हुनर तो कभी नहीं --- \omterm{def:g8:nervous-communication:synapse}{तंत्रिका-संधियाँ} हर जीवन में नए सिरे से सधती हैं। जो
जा सकता है वह है: रुझान का कच्चा माल --- \omterm{def:g1:the-five-senses:sense}{कान}, हाथ की सफ़ाई, मिज़ाज ---
यानी विरासत में मिली एक सीमा, जिसे फिर संगीत वाले घर प्रैक्टिस, तालीम और
मिसाल से भरते हैं (यानी पहुँचाने का वह दूसरा, साथ-साथ रहने वाला रास्ता)।
\end{solution}

\begin{solution}{exo:g9:heredity-and-traits:11}
पीढ़ी छोड़ना \omterm{def:g9:heredity-and-traits:heredity}{आनुवंशिकता} का सामान्य बर्ताव है: हर जनक नीली \omterm{def:g1:the-five-senses:sense}{आँख} का निर्धारक
बिना दिखाए ढो सकता है और आगे सौंप सकता है; और ऐसी दो छिपी प्रतियाँ एक ही
संतान में मिल जाएँ तो वह \omterm{def:g9:heredity-and-traits:traits}{लक्षण} दिख जाता है। यह ढर्रा बिना-दिखाए-ढोना
साबित करता है --- उस परिवार की ईमानदारी के बारे में कुछ भी नहीं।
\end{solution}

\begin{solution}{exo:g9:heredity-and-traits:12}
पिता जो कुछ पहुँचाता है वह सब \omterm{def:g8:reproductive-systems:male}{शुक्राणु} में सवार होता है, और \omterm{def:g8:reproductive-systems:male}{शुक्राणु} है
एक पूँछ और एक कसकर बँधा \omterm{def:g6:cell-unit-of-life:parts}{केंद्रक} --- इसलिए वे निर्धारक \omterm{def:g6:cell-unit-of-life:parts}{केंद्रक} में ही
बैठे होंगे। और चूँकि पूरे इंसान के पिता वाले निर्देश मिलीमीटर के
हज़ारवें हिस्सों भर के पैकेट में समा जाते हैं, इसलिए वे निर्देश अणुओं
जितनी छोटी लिखावट में लिखे होंगे --- यानी ऐसी लिपि, जो स्कूल के
\omterm{def:g5:first-look-at-cells:microscope}{सूक्ष्मदर्शी} में दिखती किसी भी कोशिका-रचना से महीन है।
\end{solution}

\begin{solution}{pb:g9:heredity-and-traits:1}
\textbf{1.} गिरजा-ठुड्डी: आनुवंशिक (परिवार का ढर्रा, और शुरू से
हाज़िर)। बाएँ हाथ का इस्तेमाल: आनुवंशिक रुझान (परिवारों के पीछे ढीले-ढाले
ढंग से चलता है, और जल्दी तय हो जाता है)। नानबाइयों की खाँसी: हासिल की हुई
(पेशे के पीछे चलती है)। रक्त-समूह: आनुवंशिक (\omterm{def:g5:flower-to-fruit:steps}{निषेचन} पर तय, और \omterm{def:g8:sexual-reproduction:gametes}{युग्मकों}
से ढोए जाने लायक़)। धूप में पके चेहरे: हासिल किए हुए (गड़ेरिये के जीवन के
पीछे चलते हैं)।

\textbf{2.} परिवार के ढर्रे वाली कसौटी: वह \omterm{def:g9:heredity-and-traits:traits}{लक्षण} किसी पेशे की लकीरों पर
बैठता है, जनन की लकीरों पर नहीं --- ब्याह कर आने वाले उसे पा लेते हैं,
जाने वाले गँवा देते हैं, और उसमें कोई \omterm{def:g8:sexual-reproduction:gametes}{युग्मक} शामिल ही नहीं।

\textbf{3.} यह कि कोई निर्धारक बिना दिखे ढोया जा सकता है: बीच वाली पीढ़ी
ने उस ठुड्डी के निर्देश आगे पहुँचाए, जबकि ख़ुद वह दूसरी ठुड्डियाँ दिखाती
रही।

\textbf{4.} हर जनक ने अपने दिखाए हुए समूह के निर्धारक के साथ-साथ एक छिपा
हुआ O वाला निर्धारक भी ढोया --- और संतान को वे दोनों छिपे हुए ही मिले।
यानी बिना दिखाए ढोना, और वह भी उसी बही की अपनी स्याही में।

\textbf{5.} वह ठुड्डी नहीं, बल्कि उसके निर्धारक --- यानी उसे बनाने के
निर्देश --- और वे भी जनन की उन इकलौती सवारियों में: हर पीढ़ी पर एक
\omterm{def:g4:animal-reproduction:sexual}{अंडाणु} और एक \omterm{def:g8:reproductive-systems:male}{शुक्राणु}, यानी सदी भर में चार हस्तांतरण।

\textbf{6.} उन चुनावों में: हर \omterm{def:g5:flower-to-fruit:steps}{निषेचन} ने साझी ठुड्डी वाले निर्धारक के
साथ-साथ बाक़ी हर चीज़ का एक ताज़ा आधा-आधा चुनाव भी बाँटा --- यानी एक ही
साझा पत्ता, और हाथ अलग-अलग।

\textbf{7.} हासिल किए हुए \omterm{def:g9:heredity-and-traits:traits}{लक्षण} कभी नहीं पहुँचते। वजह: वे शरीर के ऊतकों
में जीने से बनते हैं, और \omterm{def:g8:sexual-reproduction:gametes}{युग्मकों} --- यानी सिर्फ़ निर्धारक ढोती दो
\omterm{def:g5:first-look-at-cells:cell}{कोशिकाओं} --- के पास उन्हें दर्ज करने का कोई तरीक़ा ही नहीं।

\textbf{8.} एक ही \omterm{def:g5:flower-to-fruit:steps}{निषेचन}, और बाद में बँटवारा: एक ही निर्धारक-बँटवारा,
दो बार बना हुआ --- यानी \cref{rem:g5:human-reproduction:twins} वाले
समरूप जुड़वाँ, जिनकी गवाही बही का यह मेल दे रहा है।

\textbf{9.} \omterm{def:g8:sexual-reproduction:gametes}{युग्मक} वाला रास्ता: ठुड्डी, रक्त-समूह, और बाएँ हाथ वाला
रुझान। साथ-साथ रहने वाला रास्ता: वह खाँसी (नानबाई-ख़ाने की हवा), वे पके
चेहरे (पहाड़ों का मौसम) --- और, एलबम वाले अध्याय से, भाषाएँ, पेशे तथा
हुनर।

\textbf{10.} पीढ़ी छोड़ने से: बीच वाली पीढ़ी ने वह ठुड्डी बिना दिखाए ढोई।
और \omterm{prop:g4:journey-of-food:absorb}{रक्त} की बही से: A तथा B दिखाते जनकों ने O बिना दिखाए ढोया। दो बार,
और दोनों बार ढोना दिखाने से ज़्यादा निकला।

\textbf{11.} भौतिक रूप में: वे \omterm{def:g8:sexual-reproduction:gametes}{युग्मकों} के \omterm{def:g6:cell-unit-of-life:parts}{केंद्रकों} में बैठे हैं ---
इतना तो \omterm{def:g8:reproductive-systems:male}{शुक्राणु} की शरीर-रचना पहले ही साबित कर देती है। और वे बने किस
चीज़ से हैं, तथा इतना छोटा पैकेट इतनी बड़ी किताब लिखता कैसे है, यही तो
ठीक अगले अध्यायों का सवाल है।

\textbf{12.} ज़्यादातर \omterm{def:g9:heredity-and-traits:traits}{लक्षण} विरासत में मिली सीमाएँ हैं, जिन्हें जिया
जाता है: \omterm{def:g8:sexual-reproduction:gametes}{युग्मक} ख़ाका बाँटते हैं, और जीवन उसे भरता है।
\end{solution}
